\chapter{Conclusioni}
Il confronto tra le due architetture non ha portato a un vincitore assoluto: nel corso del capitolo sono stati analizzati vantaggi e svantaggi di entrambi gli approcci, e la scelta dipende dagli obiettivi dello sviluppatore e dal tipo di applicazione da realizzare e distribuire.

AR Foundation è la soluzione più adatta quando il team di sviluppo è ridotto, l’applicazione deve essere distribuita su più piattaforme, i tempi di produzione devono essere rapidi e si accettano prestazioni comunque buone. Se invece sono prioritari il massimo delle prestazioni, l’accesso alle funzionalità più recenti e il pieno controllo dell’hardware, con un’integrazione completa con Android o iOS, la scelta più appropriata è l’uso diretto di ARCore o ARKit. \newline

 Inoltre, l'app sviluppata si presta a diversi ambiti applicativi: l'\textbf{arredamento e l'interior design}, per visualizzare mobili e complementi nello spazio domestico prima dell'acquisto; l'\textbf{e-commerce}, per anteprime in scala reale di prodotti; l'\textbf{architettura e l'edilizia}, per la presentazione di modelli e progetti sul sito; la \textbf{didattica e la divulgazione}, per esplorare modelli di oggetti, reperti o strutture anatomiche; e l'\textbf{intrattenimento}, ad esempio in giochi o esperienze interattive in realtà aumentata.
 
 L'analisi condotta con l'\textit{Unity Profiler} \textit{on-device} mostra che l'applicazione non raggiunge il target di $60\text{ FPS}$: il tempo mediano di frame è $33.408\text{ ms}$ (circa $30\text{ FPS}$), il minimo è $19.198\text{ ms}$ e si registrano picchi isolati fino a $262.386\text{ ms}$, dovuti al rendering della camera della URP. La CPU è il limite più frequente ($41\%$ dei frame oltre il budget, contro il $20\%$ della GPU), ma il lavoro effettivo medio (\textit{Scripts} e \textit{Rendering}, circa $15.0\text{ ms}$) rientra nel budget di $16.67\text{ ms}$, mentre la quota maggiore del frame time medio è attesa di VSync ($16.039\text{ ms}$): il valore mediano è quindi in larga parte riconducibile a una cadenza di $30\text{ FPS}$, da verificare, e non solo al carico computazionale. Il carico grafico, stimato indirettamente perché il \textit{GPU Time} non è disponibile per il frame analizzato, risulta contenuto ($42$ \textit{SetPass Calls}, $7$ chiamate di disegno in $4$ batch, circa $22\text{k}$ triangoli, $46.3\text{ MB}$ di texture), così come l'occupazione di memoria ($0.66\text{ GB}$ committed). Le criticità derivano dal costo del tracciamento AR, che nel frame analizzato assorbe quasi tutto il tempo degli \texttt{Update()} ($17.07$ su $18.33\text{ ms}$), e dalle allocazioni continue sul \textit{Managed Heap}, in gran parte originate da routine di AR Foundation e del rendering, che attivano il \textit{Garbage Collector}: questo assorbe fino a $11.433\text{ ms}$ e causa buona parte dei frame oltre il budget. 
 \newline
 
 Come sviluppi futuri si individuano diverse direzioni di ottimizzazione. In
primo luogo, la verifica del frame rate target e del cap a 30~FPS: il target
è la frequenza a cui l'applicazione dovrebbe girare e determina il tempo
massimo per produrre un frame (circa 33,3~ms a 30~FPS), per cui un frame time
vicino a tale valore non indica necessariamente una CPU satura, ma può
riflettere un'attesa imposta dal limite. In secondo luogo, l'identificazione,
tramite la vista \textit{Hierarchy} del Profiler ordinata per \textit{GC
Alloc}, dell'origine delle allocazioni, cioè delle richieste di memoria sul
\textit{Managed Heap}, l'area gestita automaticamente per il codice C\#; dove
queste dipendono dal codice dell'applicazione, possono essere ridotte
mediante \textit{caching} dei risultati di operazioni costose, ovvero salvarli in una variabile i riusarli, \textit{object
pooling}, ovvero riutilizzo di oggetti attivandoli e disattivandoli al posto di
\texttt{Instantiate} e \texttt{Destroy} e riuso di \textit{buffer}, ad
esempio la lista dei risultati dei raycast.
Inoltre si può valutare l'introduzione del Garbage Collector incrementale, che distribuisce su più frame la liberazione della memoria inutilizzata, riducendo i picchi di frame time a
fronte di un lieve overhead. 
Un ulteriore intervento riguarda la riduzione del costo del tracciamento AR, ossia dell'elaborazione con cui ARCore stima la posa del dispositivo e analizza l'ambiente, ad esempio disattivando le funzionalità non necessarie come il point cloud, la stima della luce o
l'occlusione, tenendo conto del compromesso con la robustezza dell'esperienza
d'uso. 
Infine, una nuova campagna di test con build di rilascio su più dispositivi permetterebbe di ottenere misure più vicine all'esperienza reale, poiché la build di debug con Profiler collegato introduce overhead, e di quantificare i miglioramenti confrontando, prima e dopo gli interventi, frametime e memoria utilizzata.

